% Chapter wrapper — edit content in the sibling files, not here.
\chapter{Data analysis for lattice QFT}
\label{ch:data-analysis}

\todo{New chapter: jackknife/bootstrap, correlated fits, model averaging, systematic error budgets, chiral-continuum-volume extrapolation.}

\section{Anatomy of a LQCD calculation}
\begin{itemize}
\item We now have all the tools to undertake many different LQFT\aidx{lattice quantum field theory (LQFT)} calculations. Here we will walk through the steps required to perform a lattice QCD calculation of a physical quantity $\mathcal{P}$
\end{itemize}
\begin{enumerate}
    \item Define what correlation function(s) $\langle \mathcal{O} \rangle$ provides access to $\mathcal{P}$ - there may be multiple options to choose from or you may need to invent a new one. 
    \begin{itemize}
        \item It is useful to understand the precision goals of your calculation as they should influence computational choices.
    \end{itemize}
    \item Decide on what QCD theory you will use, in particular, which flavours of quark are included. Decide what discretised action (gauge \& fermion) you are going to use - ideally any action would suffice, but some choices will be better and/or more practical than others.
    \begin{itemize}
        \item e.g. if chiral symmetry\aidx{chiral symmetry} prevents operator mixing, a chiral fermion discretisation would be useful
        \item the computational costs vary significantly with this choice (otherwise every calculation would use overlap fermions\aidx{fermions!overlap})
    \end{itemize}
    \item Decide on what parameter sets you will do calculations at and how you will tune them. A full specification amounts to sets $\mathcal{S}_i = \{\text{action params, lattice geometry, statistical goal}\}$
    \begin{itemize}
        \item the key goal here is to understand what computations you will need in order to reach the continuum, infinite volume limit with physical values of the quark masses and with sufficient statistical and systematic control that your resulting $\mathcal{P}$ will be reliable.
        \item For the action parameters (bare quark masses \& gauge coupling) it is useful to define the parameter space trajectory you will explore (see Fig.~\ref{fig:ac_param_trajectories} for an example)
        \item In order to tune the quark masses to their physical values and determine the lattice spacing\aidx{lattice spacing} in physical units\aidx{scale setting}, the same number of experimentally determined quantities must be used (and computed): often the masses of various mesons and baryons are chosen along with the static quark potential\aidx{static quark potential} or the electric field energy extracted from ``Wilson flow''\aidx{Wilson flow}
        \item Ideally this tuning would be done taking the infinite volume limit at each set of action parameters
        \item This procedure is almost always undertaken iteratively, defining new parameter sets to use as more is learnt about the system.
    \end{itemize}
    \item Choose and optimise algorithms for the tasks of
    \begin{enumerate}
        \item generating representative samples of gluon configurations
        \item calculating fermion matrix inverses for construction of correlation functions\aidx{correlation function}
        \item contracting fermion inverses appropriately for the given correlation functions
    \end{enumerate}
    \begin{itemize}
        \item Different parts of the calculations have different numerical cost structures
        \item Optimised code libraries exist for some of these tasks, others must be developed
    \end{itemize}
    \item Apply for the necessary computer time/find rich benefactor and run the calculations
    \begin{itemize}
        \item Here there are also optimisations of costs and time for the computations to explore
    \end{itemize}
    \item For a fixed parameter set $\mathcal{S}_i$, perform statistical analyses to extract information about $\langle \mathcal{O} \rangle(\mathcal{S}_i)$
    \item Perform analyses of the set of $\{\langle \mathcal{O} \rangle ( \mathcal{S}_i)\}$ to arrive at the physical values of $\mathcal{P}$ in the continuum, infinite volume and physical quark mass limits
    \begin{enumerate}
        \item Statistical uncertainties at each $\mathcal{S}_i$ must be carefully propagated
        \item The forms of $a \rightarrow 0, L \rightarrow \infty, m_q \rightarrow m_q^\text{phys}$ extrapolations/interpolations can often be constrained by effective field theory methods.
        \item For different quantities (and precision goals) these extrapolations are differently challenging
    \end{enumerate}
\end{enumerate}
\begin{figure}
\redrawcompare{fig-notes-022.jpg}{fig-redraw-022.pdf}
\caption{Example: Wilson gauge and fermion action for a single quark flavour.\label{fig:ac_param_trajectories}
}
\end{figure}

\begin{itemize}
    \item The above is the case for simple, easy to define quantities but much more complicated scenarios in which the steps above do not factorise
    \begin{itemize}
        \item e.g., if $\mathcal{P}$ is determined by mass or volume dependence, analysis does not neatly factorise
    \end{itemize}
    \item Since these calculations require a lot of computing, the gauge configurations (and also in some cases the quark propagators\aidx{quark propagator}) are saved and can be reused for other physics goals - this may constrain the parameter sets $\mathcal{S}_i$ that are used
\end{itemize}

\section{State-of-the-art}

\begin{itemize}
    \item LQCD\aidx{lattice QCD (LQCD)} calculations are often carried out by large collaborations, with state of the art parameter sets for staggered/Wilson fermions\aidx{fermions!Wilson} using quark masses including the physical values and lattice geometries up to $128^3\times256$ with lattice spacing\aidx{lattice spacing}s $0.04 \text{ fm} < a < 0.15 \text{ fm}$ typical. For chiral fermions the lattice geometries are a factor of $\sim 2$ smaller in each dimension. Typically $\mathcal{O}(100)-\mathcal{O}(1000)$ somewhat independent configurations are produced
    \item Typical state of the art calculations include $N_f = 2+1$ dynamical fermions, setting $m_u = m_d \neq m_s$ and include physical values of these masses. Some calculations include charm quarks and even bottom quarks at the smallest lattice spacings. Occasionally isospin breaking is included explicitly.
    \item State-of-the-art calculations for some quantities include QED as well as QCD [since QED is not asymptotically free, it has no continuum limit\aidx{continuum limit}\aidx{triviality} and should be thought of as an effective theory]
    \item State of the art calculations use tailored variants of HMC\aidx{algorithm!hybrid Monte Carlo} for gauge field generation and typically use multigrid fermion inverters\aidx{preconditioning!multigrid}. These are implemented in community libraries

\(\left.\begin{array}{l}\text { - MILC } \\ \text { - chroma } \\ \text { - Grid + GPT (grid python toolkit) } \\ \text { - qlua } \\ \text { - openqcd }  \\ \text { - simulateQCD }  \\ \text { - tmqcd } \\ \text { - ... }\end{array}\right\}+\) quda inverters\\
These calculations use many millions of CPU/GPU hours
\end{itemize}


\section{Variance reduction and control variates}

\section{Everything is an inverse problem}
\label{sec:inverse-problems}

\todo{Section outline drafted mechanically --- prose and citations still to be
written. Delete this note once the section is yours.}

The Euclidean correlation functions a lattice calculation produces are Laplace
transforms of the spectral density they came from,
\begin{equation}
\label{eq:laplace}
  C(t) = \int_0^\infty\! dE\; \rho(E)\, e^{-E t},
\end{equation}
sampled at finitely many $t$ with correlated statistical noise. Recovering
$\rho$ from $C$ is an ill-posed inverse problem: the kernel $e^{-Et}$ maps very
different $\rho$ onto data that differ by less than the errors, so no amount of
statistics determines $\rho$ pointwise. Every extraction in this book is a case
of Eq.~\eqref{eq:laplace}; they differ only in how much of $\rho$ one needs and
therefore in how hard the problem is.

\todo{Write: this is the section's organising claim, so make it early and
concretely --- spectroscopy is not a different activity from spectral-function
reconstruction, it is the same inverse problem with a strong prior. Give the
singular-value picture of the kernel: the number of resolvable structures grows
only logarithmically with the precision of the data, which is why sub-percent
correlators still support only a handful of states.}

\subsection{Spectroscopy as the easy case}

\todo{Write: when $\rho$ is a sum of well-separated delta functions the problem
is well posed \emph{given} the prior that it is, and a multi-exponential fit is
the reconstruction. The practical machinery --- effective masses, plateau
identification, excited-state contamination, the variational method and the
GEVP --- is all inverse-problem technology under a different name. This is the
place to be honest that the prior is doing the work: the same data admit a
continuum with no isolated states, and nothing in the fit rules that out.
Cross-reference Ch.~\ref{ch:spectroscopy}, where the machinery is used, and the
finite-volume spectrum of Ch.~\ref{ch:finite-volume}, where the delta functions
are an artefact of the box.}

\subsection{Smeared spectral densities}

\todo{Write: the modern resolution --- do not ask for $\rho(E)$, ask for
$\int dE\, K_\sigma(E,\omega)\rho(E)$ with a resolution kernel $K_\sigma$ of
finite width $\sigma$, which \emph{is} determined by the data with controlled
errors, and let $\sigma \to 0$ only as the statistics allow. Backus--Gilbert
and its Hansen--Lupo--Tantalo variant, and Chebyshev approximations of the
kernel; what each assumes and what it delivers. Contrast the Bayesian route
(maximum entropy, and its successors) where the prior is explicit but the
result depends on it.}

\subsection{Where the hard cases arise}

\todo{Write: a short catalogue, each one a pointer rather than a treatment.
\begin{itemize}
  \item \emph{Transport coefficients} --- shear and bulk viscosity from the
    small-$\omega$ limit of $\rho$, the hardest case of all
    (Ch.~\ref{ch:thermodynamics}).
  \item \emph{Inclusive decay rates} --- a phase-space integral of $\rho$
    against a smooth kernel, so a smeared density suffices; this is why
    inclusive semileptonic decays became tractable
    (Ch.~\ref{ch:tests-of-sm}).
  \item \emph{Hadronic vacuum polarisation} --- an integral of the correlator
    against a known weight, which is the \emph{easy} direction of the same
    transform, and the reason window observables work.
  \item \emph{Finite-volume scattering} --- L\"uscher's method inverts a
    different map, from the finite-volume spectrum to the phase shift, but the
    conditioning question is the same (Ch.~\ref{ch:finite-volume}).
\end{itemize}
The point to land: the difficulty is set by how smooth a functional of $\rho$
is wanted, not by the quantity's physical importance.}

\subsection{Diagnostics}

\todo{Write: how to tell whether an answer is data-driven or prior-driven.
Vary the kernel width and watch the error grow; check stability against the fit
range and the number of states; use synthetic data with a known $\rho$ and the
covariance of the real ensemble. Connect to model averaging earlier in this
chapter --- an information criterion over fit ranges is a way of admitting that
the prior is uncertain, not a way of removing it.}